\documentclass[11pt,a4paper]{article}
\usepackage[T1]{fontenc}
\usepackage[utf8]{inputenc}
\usepackage{amsmath,amssymb,xcolor}
\usepackage[margin=25mm]{geometry}
\usepackage[colorlinks=true,urlcolor=blue]{hyperref}
\definecolor{nexorablue}{HTML}{2459B5}
\setlength{\parindent}{0pt}
\setlength{\parskip}{7pt}
\setlength{\emergencystretch}{2em}
\renewcommand{\familydefault}{\sfdefault}
\begin{document}

{\small\bfseries\color{nexorablue} NEXORA CPGE}\par
{\small MPSI / NEX-RGT-04 / Corrigé}\par
{\LARGE\bfseries Une bobine en deux phases}\par
\vspace{0.25em}\hrule\vspace{0.6em}

\section*{Charge}
Le courant limite vaut $2\,\mathrm A$. À la commutation, $i_1=2(1-e^{-2})\simeq1{,}73\,\mathrm A$. Le courant est continu à cet instant.

\section*{Décharge et énergie}
Dans la seconde phase, $\tau_2=0{,}12/9\simeq13{,}3\,\mathrm{ms}$ et $i=i_1e^{-\Delta t/\tau_2}$. Le seuil est atteint après $\Delta t=\tau_2\ln(i_1/0{,}25)\simeq25{,}8\,\mathrm{ms}$. L’énergie dissipée jusqu’à la fin vaut $Li_1^2/2\simeq0{,}179\,\mathrm J$.


\end{document}
